\section{自适应 RAG 的三类关键决策}
\label{sec:adaptive}

\subsection{何时检索：路由、动态触发与反思}

\method{输入路由。} Adaptive-RAG 训练较小的语言模型预测问题复杂度，在不检索、单次检索和多步迭代之间选择路径~\citep{jeong2024adaptive}。它适合请求类型较稳定的问答服务，但句式复杂度与知识缺口并不总是一致：短问题可能依赖实时信息，包含多个实体的问题也可能都在模型熟悉的范围内。路由标签因此不是问题文本的固有属性：模型参数知识、语料覆盖和引用要求都会改变最优路径，更换这些条件后需要重新校准。

\method{生成中触发。} FLARE 先产生临时下一句，出现低概率词元时用其检索并重新生成~\citep{jiang2023flare}。DRAGIN 将信息需求检测与查询构建分开，前者结合词元不确定性、语义和自注意力，后者从完整上下文中选择重要内容形成查询~\citep{su2024dragin}。二者适合写作过程中逐渐出现的信息需求。概率信号也会受到措辞和解码方式影响，低概率未必意味着事实缺失，高概率同样可能对应错误事实。

\method{学习检索与反思。} Self-RAG 用特殊反思标记控制按需检索，并评价文档相关性、答案支持度和输出有用性~\citep{asai2024selfrag}。它将证据评价延伸到生成片段，代价是反思数据构造与额外训练。三种方式分别在输入、生成过程和模型训练层面实现控制，选择取决于请求分布、文本长度和可用训练资源。

\subsection{检索什么：查询表示与结构修正}

HyDE 先生成假设性回答文档，再用无监督向量编码器检索相近的真实文档~\citep{gao2023hyde}。假设内容可能有误，却能提供目标答案附近的术语和语义模式，最终回答仍依赖真实文档。RQ-RAG 将改写、问题分解、歧义消解和按需搜索统一为带控制标记的生成过程~\citep{chan2024rqrag}。前者主要改变检索表示，后者还改变问题结构，可处理对话指代和相互依赖的子问题。

改写可能擅自补充时间或实体，也可能删去限制条件。评价查询模块时，除检索相关性外，还应检查原始约束是否保留，并记录每轮新增的假设，以定位检索偏差的来源。

\subsection{证据是否可用：评价、过滤与纠正}

CRAG 在检索器和生成器之间加入轻量评价器，将结果分为 Correct、Incorrect 和 Ambiguous，分别执行知识精炼、替换搜索和混合处理；分解—重组步骤用于筛选关键信息~\citep{yan2024crag}。检索失败由此对应具体的恢复行动。

相关性和充分性需要分开。多篇文档可能重复同一事实，留下另一个子问题没有答案；增加文档又会加重噪声和位置偏差~\citep{liu2024lost,cuconasu2024noise}。压缩和重排应保留支持关键事实的片段及来源，冲突文档则需记录分歧。

RAG-Critic 建立层次化错误体系，用由粗到细训练的评价模型生成反馈，再由规划与执行模块选择纠正流程~\citep{dong2025ragcritic}。它与 CRAG、Self-RAG 的差异在于反馈组织方式：检索质量、答案支持度和细粒度错误分别对应不同检查。评价器进入控制回路后，误判也会直接改变后续行动，调用成本与评价可靠性需要一并考虑。

S2G-RAG 和 FAIR-RAG 进一步将证据评价写成结构化状态：前者输出信息缺口并据此形成查询，后者把答案需求拆成检查项，区分已确认事实与待补材料~\citep{li2026s2g,aghajani2025fairrag}。这类状态使“继续找什么”可以接受检查，但内容覆盖只是充分性的一个维度。要求可核验引用时，还需记录支持片段与来源；处理动态知识时，日期和来源冲突也会改变证据能否用于作答。充分性判断因而依赖任务要求，不能仅由相关文档数量或槽位填充比例决定。

\subsection{代表方法的统一比较}

表~\ref{tab:methods} 比较上述方法的控制对象与使用代价，搜索策略学习和停止控制在第~\ref{sec:iterative} 节继续讨论。

\begin{table}[t]
\centering
\small
\caption{典型自适应 RAG 方法的控制机制比较}
\label{tab:methods}
\begin{tabularx}{\textwidth}{p{2.4cm}p{2.15cm}XX}
\toprule
方法 & 主要决策 & 核心机制与适用场景 & 主要使用代价或失败情形 \\
\midrule
Adaptive-RAG & 输入路径选择 & 按预测复杂度路由到无检索、单次或迭代流程 & 跨模型、跨语料部署时需重新校准 \\
FLARE & 生成中触发 & 低概率词元触发检索，以临时下一句形成查询 & 表达变化可能引起额外检索 \\
DRAGIN & 触发与查询 & 综合不确定性、语义和注意力识别需求并构建查询 & 控制信号较多，实现与调参成本较高 \\
Self-RAG & 检索、证据与输出 & 学习反思标记，评价相关性、支持度和有用性 & 依赖反思数据构造和额外训练 \\
HyDE & 查询表示 & 用假设文档连接问题与真实文档的向量空间 & 错误假设可能将检索引向错误主题 \\
RQ-RAG & 查询修正 & 统一改写、分解、消歧与按需搜索 & 改写及轨迹选择错误会影响后续证据 \\
CRAG & 证据评价与恢复 & 根据检索质量执行精炼、替换或补充搜索 & 评价误判会改变恢复路径 \\
IRCoT & 迭代查询 & 交替执行推理和检索，让中间结果引导查询 & 多轮调用成本较高，中间错误可能传播 \\
\bottomrule
\end{tabularx}
\end{table}
